\section{Fórmula de Bayes y descomposición del score condicional}

La herramienta matemática central para resolver problemas inversos es la \textbf{fórmula de Bayes}, que descompone el score condicional en la suma del score marginal y el score de verosimilitud.

\subsection{Descomposición bayesiana del score condicional}

Nuestro objetivo es calcular el score condicional $\nabla_{x_t} \log p(x_t | y)$ para guiar el proceso de generación durante la eliminación de ruido. Utilizando la fórmula de Bayes:

$$
p(x_t | y) = \frac{p(y | x_t) \cdot p(x_t)}{p(y)}
$$

Tomando logaritmos en ambos lados y calculando el gradiente:

$$
\nabla_{x_t} \log p(x_t | y) = \underbrace{\nabla_{x_t} \log p(x_t)}_{\text{Score marginal (conocido)}} + \underbrace{\nabla_{x_t} \log p(y | x_t)}_{\text{Score de verosimilitud (por calcular)}}
$$

\begin{importantbox}{Score condicional = Score marginal + Score de verosimilitud}
\begin{itemize}
    \item \textbf{Score marginal} $\nabla_{x_t} \log p(x_t)$: esta es exactamente la función de score que los modelos de difusión preentrenados ya aprendieron a predecir y se puede obtener directamente del modelo.
    \item \textbf{Score de verosimilitud} $\nabla_{x_t} \log p(y | x_t)$: mide qué tan probable es que la observación $y$ aparezca dado el estado intermedio actual $x_t$. Esta es la parte que requiere un cálculo adicional (o aproximado).
\end{itemize}
Por lo tanto, todo el problema se reduce a: \textbf{¿cómo aproximar el score de verosimilitud sin entrenar una red neuronal adicional?}
\end{importantbox}

\subsection{Comparación con Classifier Guidance}

\begin{knowledgebox}{Métodos existentes bajo un marco unificado}
Las técnicas de guiado estudiadas anteriormente pueden entenderse de manera unificada bajo este marco:
\begin{itemize}
    \item \textbf{Classifier Guidance}: estimando directamente el score de verosimilitud $\nabla_{x_t} \log p(y | x_t)$ entrenando una red neuronal de clasificación adicional.
    \item \textbf{Classifier-Free Guidance}: aprendiendo conjuntamente con la misma red tanto el ruido condicional como el incondicional, calculando implícitamente el score de verosimilitud a través de la diferencia entre ambos.
\end{itemize}
Ambos métodos requieren entrenamiento. Lo que nosotros vamos a hacer es \textbf{aproximar el score de verosimilitud sin necesidad de entrenamiento}.
\end{knowledgebox}

\subsection{Resumen de la sección}

La descomposición bayesiana nos indica que la generación condicional solo requiere añadir un score de verosimilitud adicional sobre el score marginal del modelo preentrenado. El desafío central restante es: ¿cómo aproximar este score de verosimilitud de manera eficiente?



